\chapter{Perceptual Reconstruction and Harmonic Completion}
\label{app:perceptual}

\epigraph{A signal may remain semantically recognizable while losing fine-grained
structural relations necessary for perceptual richness, temporal immersion, or
sustained attentional coherence.}{}

\noindent Perception does not operate through perfect recovery of external
structure. All reconstructive systems operate under informational incompleteness,
energetic limitation, and finite temporal accessibility. Perceptual systems
therefore rely upon partial reconstruction procedures that stabilize coherent
interpretations from degraded or underspecified inputs. This appendix investigates
harmonic reconstruction, perceptual coherence, and temporal continuity under
compressed representation regimes.

The material here should be understood as \textbf{[P]} phenomenological
interpretation and \textbf{[H]} heuristic engineering throughout. It does not
establish that compressed audio damages consciousness, removes emotional essence,
or eliminates hidden biological frequencies necessary for cognition. Such claims
exceed available evidence. The framework proposes only that reconstruction quality
influences perceptual coherence and that harmonic continuity may affect subjective
richness, immersion, and attentional stability under constrained representational
conditions.

\section{Compression and the Symbolic/Inferential Distinction}

All lossy compression systems preserve certain structural relations while discarding
others. The resulting reconstruction is shaped not only by retained information but
also by the assumptions embedded within the reconstruction procedure itself.

Perceptual degradation frequently emerges when temporally or harmonically significant
relations are removed despite preserving coarse symbolic intelligibility. This
reflects the broader separation between symbolic continuity and inferential
continuity developed in Chapter~\ref{ch:scaffolded}. A signal may remain
semantically recognizable---the words are audible, the melody is identifiable---
while losing fine-grained structural relations necessary for perceptual richness,
temporal immersion, or sustained attentional coherence.

\section{Harmonic Reconstruction as Constraint Completion}

Suppose a signal admits a partial harmonic decomposition:
\[
  x(t) = \sum_{k=1}^{N} a_k \cos(\omega_k t + \phi_k).
\]
Compression removes a subset of components:
\[
  \hat{x}(t) = \sum_{k \in S} a_k \cos(\omega_k t + \phi_k), \quad S \subseteq \{1,\ldots,N\}.
\]
Reconstruction seeks admissible completion operators $\mathcal{R}: \hat{x}(t) \mapsto
\tilde{x}(t)$ that infer missing relational structure from surviving harmonic
constraints:
\[
  \tilde{x} = \arg\min_{y \in \mathcal{C}} D(y, \hat{x}),
\]
where $\mathcal{C}$ denotes the space of temporally coherent reconstructions and
$D$ extracts perceptually relevant relational structure rather than raw amplitude
similarity alone.

The perceptual divergence between original and reconstruction is:
\[
  D(x, \tilde{x}) = \int \|K(x(t)) - K(\tilde{x}(t))\|^2\, dt,
\]
where $K$ extracts perceptually relevant relational structure. This is a
constrained interpolation problem, not a restoration problem. The reconstructed
signal is an inferential approximation generated under bounded informational
accessibility, not a recovery of the original.

\section{Temporal Continuity and Perceptual Stability}

Biological perception exhibits strong sensitivity to temporal discontinuity.
Abrupt fragmentation, phase instability, and incoherent transient structure
increase perceptual fatigue and reduce sustained attentional stabilization.
Define local temporal fragmentation through phase structure $\phi(t)$:
\[
  F(t) = \left|\frac{d^2\phi}{dt^2}\right|.
\]
Perceptual stability emerges when integrated fragmentation remains bounded:
\[
  \int_{t_0}^{t_1} F(t)\,dt < \epsilon.
\]

\claimH Harmonic reconstruction mechanisms may partially restore such continuity
even when substantial information has been lost. The resulting increase in perceived
richness does not imply that hidden information has been recovered in a literal
sense. Rather, the reconstruction process stabilizes relational structures
compatible with prior perceptual expectations and surviving signal constraints.

\section{Inferential Completion in Auditory Perception}

Perceptual reconstruction depends heavily upon inferential completion. Human
auditory systems routinely interpolate missing information from contextual
structure, harmonic expectation, and temporal continuity:
\[
  \hat{x} = \arg\max_z P(z \mid y, H),
\]
where $y \subset x$ is the available perceptual input and $H$ represents prior
structural expectations. Hallucination, interpolation, and perceptual completion
all arise naturally whenever $\dim(y) < \dim(x)$. The system operates by
constrained reconstruction rather than exhaustive recovery of inaccessible total
states. This is exactly the reconstruction bound of Chapter~\ref{ch:vocabulary}
applied to the auditory domain.

\claimP Such effects belong primarily to psychoacoustics, perceptual inference,
and temporal coherence theory. They do not require speculative metaphysics
about consciousness or hidden frequencies. The interesting questions are
computational and psychoacoustic: which structural relations does the auditory
system most depend on for stability, and which reconstruction heuristics best
preserve those relations under compression?
